\chapter{เทคโนโลยีการจำลองเสมือนจริง AR/XR และ WebXR ในการสอนฟิสิกส์ (AR/XR Frameworks for Physics)}
\label{chap:ar_xr_frameworks}
\index{WebXR}
\index{Three.js}
\index{MediaPipe}
\index{การจำลองเสมือนจริง}

\begin{chapterobjectives}
เมื่อศึกษาบทที่ 4 จบแล้ว นิสิตและผู้ศึกษาสามารถ:
\begin{enumerate}
    \item อธิบายสถาปัตยกรรมของ WebXR Device API และวงจรชีวิตของเซสชัน AR/VR
    \item ประยุกต์ใช้ไลบรารี Three.js ในการสร้างฉากสามมิติ เรนเดอร์โมเดลฟิสิกส์ และควบคุมแสงเงาแบบ PBR (Physically Based Rendering)
    \item ผสานระบบตรวจจับข้อต่อนิ้วมือ 21 จุด 3 มิติ (MediaPipe Hands) สำหรับการโต้ตอบกับแบบจำลองฟิสิกส์แบบไร้การสัมผัส
    \item ออกแบบวงรอบแอนิเมชันความเร็ว 60 เฟรมต่อวินาที ที่ไม่มีขยะในหน่วยความจำ (Zero-Garbage Collection Loop)
    \item พัฒนาเว็บแอปพลิเคชัน AR แบบ Standalone เพื่อจำลองการเปลี่ยนระดับพลังงานในอะตอมของบอร์ (Bohr Atom Simulator)
\end{enumerate}
\end{chapterobjectives}

\section{วิวัฒนาการและบทบาทของ WebXR ในการศึกษาฟิสิกส์}
\index{Spatial Computing}
\index{ความเป็นจริงเสริม}

มโนทัศน์ในวิชาฟิสิกส์แผนใหม่ เช่น ปริภูมิ-เวลา 4 มิติ ความน่าจะเป็นของออร์บิทัลอิเล็กตรอน หรือการสลายตัวของนิวเคลียส เป็นปรากฏการณ์ที่ไม่สามารถมองเห็นได้ด้วยตาเปล่าในชีวิตประจำวัน การใช้เทคโนโลยี \textbf{ความเป็นจริงขยาย (Extended Reality: XR)} ซึ่งครอบคลุมทั้ง Augmented Reality (AR) และ Virtual Reality (VR) จึงเปิดมิติใหม่ในการแปลงนามธรรมทางคณิตศาสตร์ให้กลายเป็นวัตถุเชิงพื้นที่ (Spatial Objects) ที่ผู้เรียนสามารถสังเกตและมีปฏิสัมพันธ์ได้รอบทิศทาง

\textbf{WebXR Device API} คือมาตรฐานเว็บสากลของ W3C ที่อนุญาตให้เว็บเบราว์เซอร์สามารถเข้าถึงเซนเซอร์ตรวจจับตำแหน่งและการเคลื่อนไหว 6 แกนอิสระ (6-DoF Tracking) ของสมาร์ตโฟนหรือแว่นตาอัจฉริยะ (Head-Mounted Displays) ได้โดยตรง โดยผู้ใช้ไม่ต้องติดตั้งแอปพลิเคชันเสริมใดๆ

\begin{figure}[H]
\centering
\begin{tikzpicture}[scale=0.95]
    % Real World Layer
    \draw[darkNavy, fill=gray!10, rounded corners=6pt] (-5, -1.8) rectangle (5, -0.6);
    \node at (0, -1.2) [font=\footnotesize\bfseries\color{darkNavy}] {สภาพแวดล้อมโลกจริง (Physical World \& Camera Video Feed)};
    
    % WebXR Hardware Abstraction Layer
    \draw[cyberBlue, fill=cyan!10, rounded corners=6pt] (-5, 0.0) rectangle (5, 1.2);
    \node at (0, 0.6) [font=\footnotesize\bfseries\color{cyberBlue}] {WebXR Device API (Pose, Hit-Test, View Matrix, Projection Matrix)};
    
    % Three.js & Shader Engine
    \draw[neonPurple, fill=purple!10, rounded corners=6pt] (-5, 1.8) rectangle (5, 3.0);
    \node at (0, 2.4) [font=\footnotesize\bfseries\color{neonPurple}] {Three.js Scene Graph \& Custom GLSL Physics Shaders};
    
    % Interaction Layer (MediaPipe)
    \draw[emeraldGlow, fill=green!10, rounded corners=6pt] (-5, 3.6) rectangle (5, 4.8);
    \node at (0, 4.2) [font=\footnotesize\bfseries\color{emeraldGlow!80!black}] {MediaPipe Hands 21-Joint 3D Gesture Controller (Zero-Touch Interface)};
    
    % Connecting arrows
    \draw[-{Stealth[scale=1.2]}, line width=1.5pt, slateGrey] (0, -0.6) -- (0, 0.0);
    \draw[-{Stealth[scale=1.2]}, line width=1.5pt, slateGrey] (0, 1.2) -- (0, 1.8);
    \draw[-{Stealth[scale=1.2]}, line width=1.5pt, slateGrey] (0, 3.0) -- (0, 3.6);
\end{tikzpicture}
\caption{สถาปัตยกรรมซอฟต์แวร์จำลองเสมือนจริง WebXR ร่วมกับ Three.js และ MediaPipe Hands}
\label{fig:webxr_architecture}
\end{figure}

\section{วงจรชีวิตของ WebXR Session และการเรนเดอร์ 3 มิติ}
\index{WebXR Session}
\index{Hit-Test}

ขั้นตอนการทำงานของแอปพลิเคชัน WebXR มีลำดับดังต่อไปนี้:
\begin{enumerate}
    \item \textbf{ตรวจสอบการรองรับอุปกรณ์ (Device Capability Check):} เรียกใช้ฟังก์ชัน \\
    \texttt{navigator.xr.isSessionSupported('immersive-ar')}
    \item \textbf{ร้องขอเปิดเซสชัน (Request Session):} เรียกใช้คำสั่ง \\
    \texttt{navigator.xr.requestSession('immersive-ar', \{requiredFeatures: ['hit-test']\})}
    \item \textbf{การจับพื้นผิวจริง (Hit-Testing):} ยิงลำแสงเสมือน (Virtual Raycast) จากกึ่งกลางหน้าจอไปกระทบพื้นผิวจริง เพื่อกำหนดพิกัดอ้างอิงของจุดศูนย์กลางการทดลอง (Anchor Pose)
    \item \textbf{วงรอบการเรนเดอร์ (Render Loop):} รับค่าการขยับของกล้องผ่าน \\
    \texttt{XRSession.requestAnimationFrame(onXRFrame)} เพื่อปรับปรุงการเรนเดอร์ของ Three.js ให้ตรงกับมุมมองของผู้ใช้อย่างแม่นยำ
\end{enumerate}

\section{การควบคุมแบบไร้สัมผัสด้วย MediaPipe Hands}
\index{MediaPipe}
\index{Hand Tracking}

MediaPipe Hands ของ Google ใช้โครงข่ายประสาทเทียมแบบ Deep Learning ที่ได้รับการเทรนให้ตรวจจับจุดข้อต่อนิ้วมือ 21 จุดในพิกัด 3 มิติ ($x, y, z$) จากเฟรมวิดีโอกล้องหน้าหรือหลังแบบเรียลไทม์
ท่าทางสำคัญที่นำมาใช้ในการควบคุมการทดลองฟิสิกส์:
\begin{itemize}
    \item \textbf{Pinch Gesture (นิ้วโป้งแตะปลายนิ้วชี้):} ใช้จับอนุภาค ยืดบ่อศักย์ หรือหมุนกรวยแสง
    \item \textbf{Palm Facing (หงายฝ่ามือ):} ใช้เป็นแท่นรองรับอนุภาคเสมือนหรือกล่องพลังงาน
    \item \textbf{Index Pointing (ชี้ตำแหน่ง):} ใช้เลือกสถานะควอนตัม ($n=1, 2, 3$)
\end{itemize}

\section{วิศวกรรมลูปประสิทธิภาพสูง (Zero-GC Loop Engineering)}
\index{Zero-GC}
\index{Garbage Collection}

ในแอนิเมชัน 60 FPS เบราว์เซอร์มีเวลาประมวลผลเพียง $16.67\text{ ms}$ ต่อเฟรม หากโค้ดสร้างอ็อบเจกต์ใหม่ในลูป เช่น \texttt{new THREE.Vector3()} ขยะในหน่วยความจำจะสะสมอย่างรวดเร็วและกระตุ้นการทำงานของ Garbage Collector (GC Pause) ส่งผลให้ภาพกระตุกอย่างรุนแรง

\begin{pitfallbox}[การจัดสรรหน่วยความจำในลูปแอนิเมชัน 60 FPS]
\textbf{ข้อพึงปฏิบัติอย่างเคร่งครัด:}
\begin{itemize}
    \item ประกาศตัวแปรชั่วคราว เช่น เวกเตอร์และเมทริกซ์ ไว้นอกลูปหลักเสมอ (Global/Module Scope Scratchpad)
    \item ใช้เมธอด \texttt{.copy()}, \texttt{.set()}, หรือ \texttt{.add()} แทนการสร้างอินสแตนซ์ใหม่
    \item ลบ Geometry และ Material ที่ไม่ได้ใช้ออกจาก VRAM ด้วยคำสั่ง \texttt{.dispose()}
\end{itemize}
\end{pitfallbox}

\section{ตัวอย่างโค้ดระบบเต็ม: แบบจำลองอะตอมบอร์เสมือนจริงใน WebXR}
\index{Bohr Atom AR}
\index{JavaScript}

\begin{pythonbox}[bohr\_atom\_webxr.html: โครงสร้างโค้ดจำลองอะตอมบอร์แบบสมบูรณ์]
<!DOCTYPE html>
<html lang="th">
<head>
    <meta charset="UTF-8">
    <title>WebXR Bohr Atom Simulator - Dr. Chewa Thassana</title>
    <script src="https://cdnjs.cloudflare.com/ajax/libs/three.js/r128/three.min.js"></script>
    <style>
        body { margin: 0; overflow: hidden; font-family: 'Sarabun', sans-serif; }
        #ar-button {
            position: absolute; bottom: 20px; left: 50%;
            transform: translateX(-50%); padding: 12px 24px;
            background: #0284C7; color: white; border: none;
            border-radius: 8px; font-weight: bold; font-size: 16px;
        }
    </style>
</head>
<body>
    <button id="ar-button">เข้าสู่โหมดการทดลอง AR</button>
    <script>
        let scene, camera, renderer, xrSession = null;
        let electron, nucleus, orbitLines = [];
        let currentLevel = 1, targetRadius = 1.0;
        const scratchVec = new THREE.Vector3(); // Zero-GC scratchpad

        function initScene() {
            scene = new THREE.Scene();
            camera = new THREE.PerspectiveCamera(70, window.innerWidth/window.innerHeight, 0.01, 20);
            
            renderer = new THREE.WebGLRenderer({ antialias: true, alpha: true });
            renderer.setSize(window.innerWidth, window.innerHeight);
            renderer.xr.enabled = true;
            document.body.appendChild(renderer.domElement);

            // นิวเคลียส (โปรตอน)
            const nucGeo = new THREE.SphereGeometry(0.06, 32, 32);
            const nucMat = new THREE.MeshStandardMaterial({ color: 0xF59E0B, roughness: 0.2 });
            nucleus = new THREE.Mesh(nucGeo, nucMat);
            nucleus.position.set(0, 0, -1.0);
            scene.add(nucleus);

            // อิเล็กตรอน
            const elecGeo = new THREE.SphereGeometry(0.02, 16, 16);
            const elecMat = new THREE.MeshBasicMaterial({ color: 0x0284C7 });
            electron = new THREE.Mesh(elecGeo, elecMat);
            scene.add(electron);

            // แสงสว่าง
            const light = new THREE.PointLight(0xffffff, 1.5, 10);
            light.position.set(1, 2, 0);
            scene.add(light);
            scene.add(new THREE.AmbientLight(0x404040));
        }

        let angle = 0;
        function renderLoop(timestamp, frame) {
            angle += 0.03;
            const r = 0.3 * (currentLevel ** 2);
            electron.position.x = nucleus.position.x + r * Math.cos(angle);
            electron.position.z = nucleus.position.z + r * Math.sin(angle);
            electron.position.y = nucleus.position.y;

            renderer.render(scene, camera);
        }

        document.getElementById('ar-button').addEventListener('click', async () => {
            if (navigator.xr) {
                const session = await navigator.xr.requestSession('immersive-ar', {
                    requiredFeatures: ['local-floor']
                });
                renderer.xr.setSession(session);
                renderer.setAnimationLoop(renderLoop);
            } else {
                alert('อุปกรณ์ไม่รองรับ WebXR กรุณาใช้สมาร์ตโฟนที่รองรับ ARCore');
            }
        });

        initScene();
    </script>
</body>
</html>
\end{pythonbox}

\section{ตัวอย่างการแก้ปัญหาเชิงเทคนิคตามระเบียบวิธี P-S-C-T}

\begin{psctexample}[การปรับค่าพิกัดการแปลงพื้นที่สำหรับ AR Tracking ให้เรียบเนียน (Smooth Interpolation)]
\textbf{โจทย์ (Problem):} \\
ในการจำลอง AR เมื่อเครื่องมือตรวจวัด (เช่น MediaPipe หรือ WebXR Hit-Test) ส่งตำแหน่งจุดสัมผัส ($x_t, y_t, z_t$) เข้ามาทุกเฟรม มักพบปัญหาการสั่นไหวเล็กน้อยของข้อมูล (Jittering / Noise) ส่งผลให้อิเล็กตรอนเสมือนสั่นกระตุกจนสูญเสียความสมจริง จงออกแบบตัวกรองทางคณิตศาสตร์ที่คำนวณได้อย่างรวดเร็วในรอบเฟรม 60 FPS เพื่อให้การเคลื่อนที่ราบรื่น

\textbf{ยุทธวิธี (Strategy):} \\
ใช้ตัวกรองแบบชะลอชี้กำลังหรือการประมาณค่าในช่วงเชิงเส้น (Linear Interpolation: LERP) ระหว่างตำแหน่งปัจจุบัน $\mathbf{P}_{\text{curr}}$ และตำแหน่งเป้าหมาย $\mathbf{P}_{\text{target}}$:
\begin{equation}
\mathbf{P}(t + \Delta t) = \mathbf{P}(t) + \alpha \left( \mathbf{P}_{\text{target}} - \mathbf{P}(t) \right) = (1 - \alpha)\mathbf{P}(t) + \alpha \mathbf{P}_{\text{target}}
\end{equation}
โดยที่ $\alpha \in (0, 1]$ คือแฟกเตอร์การผ่อนคลาย (Smoothing Factor)

\textbf{การคำนวณ (Computation):} \\
หากต้องการให้การตอบสนองมีความหน่วง 95\% เข้าสู่เป้าหมายภายในเวลา $T = 0.2\text{ วินาที}$ ที่ความถี่เฟรม $60\text{ Hz}$ ($\Delta t = 1/60\text{ s}$ หรือมีจำนวน 12 เฟรม):
\begin{equation}
(1 - \alpha)^{12} \le 0.05 \implies 1 - \alpha = (0.05)^{1/12} \approx 0.779 \implies \alpha \approx 0.22
\end{equation}
ดังนั้น การเขียนโค้ด LERP ในลูปแอนิเมชันคือ:
\begin{equation}
\mathbf{P}_{\text{curr}}.\text{lerp}(\mathbf{P}_{\text{target}}, 0.22)
\end{equation}

\textbf{ข้อคิดทางฟิสิกส์และวิศวกรรม (Takeaway):} \\
การคำนวณ LERP แบบเวกเตอร์ใช้การบวกและคูณสเกลาร์เพียง 3 มิติ ซึ่งใช้เวลาน้อยกว่า $0.001\text{ ms}$ บน CPU จึงไม่กระทบต่อ Frame Budget ของเบราว์เซอร์ แต่สามารถกำจัดความสั่นไหวของเซนเซอร์ AR ได้อย่างมีประสิทธิภาพสูงสุด
\qedexam
\end{psctexample}

\section{สรุปสาระสำคัญประจำบท}

\begin{table}[H]
\centering
\begin{tabularx}{\textwidth}{l Z Y}
\toprule
\textbf{เทคโนโลยี / ฟังก์ชัน} & \textbf{เครื่องมือหลัก} & \textbf{หน้าที่ในการศึกษาฟิสิกส์} \\
\midrule
WebXR Device API & W3C Standard & เชื่อมต่อฮาร์ดแวร์ AR/VR บนเบราว์เซอร์โดยตรง \\
Three.js & WebGL / WebGPU & จัดการฉาก เรนเดอร์วัสดุ แสงเงา และตาข่ายสามมิติ \\
MediaPipe Hands & Machine Learning & ติดตามข้อต่อนิ้วมือ 21 จุด เพื่อควบคุมการทดลองไร้สัมผัส \\
Zero-GC Loop & Object Pooling & ป้องกันการกระตุกของภาพ รักษาอัตราเฟรมคงที่ 60 FPS \\
Spatial LERP & $(1-\alpha)\mathbf{x}_0 + \alpha\mathbf{x}_1$ & ลดสัญญาณรบกวนจากการตรวจจับตำแหน่งของกล้อง AR \\
\bottomrule
\end{tabularx}
\caption{สรุปเครื่องมือและเทคนิคหลักในการสร้างแบบจำลองฟิสิกส์ด้วย WebXR}
\label{tab:ar_xr_summary}
\end{table}

\section{แบบฝึกหัดท้ายบท}
\begin{enumerate}
    \item จงอธิบายความแตกต่างระหว่าง AR แบบมาร์กเกอร์ (Marker-based AR) และ AR แบบตรวจจับระนาบพื้นผิว (Surface / Markerless AR via Hit-Testing) ในแง่ของการนำมาประยุกต์ใช้ในการทดลองฟิสิกส์
    \item หากต้องการจำลองสนามแม่เหล็ก $B$ รอบลวดตัวนำที่มีกระแสไฟฟ้าไหลผ่าน จงร่างอัลกอริทึมในการแสดงเส้นแรงแม่เหล็กเสมือนจริงในมุมมอง 3 มิติ
    \item เหตุใดการสร้างอินสแตนซ์เวกเตอร์ใหม่ (\texttt{new THREE.Vector3()}) ภายในฟังก์ชัน \texttt{requestAnimationFrame} จึงก่อให้เกิดปัญหาต่อประสบการณ์ของผู้ใช้ และมีแนวทางแก้ไขอย่างไร?
\end{enumerate}
